\section{Full first-order iPRA and both interpretation composites}
\label{sec:interpretations}

The comparison is with full first-order intuitionistic primitive recursive
arithmetic, not just its equational fragment or its closed numerical
consequences.  This section gives the interpretation maps and identifies
the inherited mathematical work needed for the reverse direction.

\subsection{The target theory}

Use ordinary intuitionistic first-order logic with actual equality and
falsity.  Its language has $0,S$ and a symbol for every fixed primitive
recursive definition, generated from zero, successor, projections,
composition, and primitive recursion.  The nonlogical axioms are those
definitions' equations, successor separation and injectivity,
\[
   Sx=0\to\bot,\qquad Sx=Sy\to x=y,
\]
and induction for every quantifier-free formula $A(n,\vec p)$:
\[
 \bigl(A(0,\vec p)\wedge
       \forall n(A(n,\vec p)\to A(Sn,\vec p))\bigr)
       \to\forall n A(n,\vec p).
\]
The full logic admits the usual introduction and elimination rules for
$\wedge,\vee,\to,\forall,\exists$, falsity elimination, assumption use,
equality reflexivity, and equality substitution.  Universal introduction
and existential elimination have their ordinary fresh-variable conditions.
Excluded middle for arbitrary formulas is not a rule.  This is the theory
deeply embedded as \texttt{PR} in the pinned repository.

\subsection{An intermediate ordinary scalar theory}

Let $T_I$ be intuitionistic first-order equality logic on the scalar
iterator terms, with the two computation equations, coverage,
\[
             \forall x\,(x=0\vee\exists y\,x=Sy),
\]
separation $(0=S0)\to\bot$, and contextual induction for literal
equations.  Elaborate a fixed update term $h(u,\vec p)$ into an ordinary
function symbol $i_h(n,a,\vec p)$ when an orthodox first-order signature is
desired.  Its defining equations are
\[
 i_h(0,a,\vec p)=a,\qquad
 i_h(Sn,a,\vec p)=h(i_h(n,a,\vec p),\vec p).
\]
These declarations are staged by finite syntax.  There is no numerical
sort of arbitrary update programs and no universal evaluator.

The interpretation of native syntax into $T_I$ leaves numerical terms and
equations unchanged and sends
\[
 D_{x,y}B\quad\mapsto\quad
 \exists x\bigl(B^\circ(x,0)\to\forall y B^\circ(x,y)\bigr).
\]
An unconditional clause becomes $\forall y B^\circ(y)$; a conditional
clause becomes $A^\circ\to\forall y B^\circ(y)$.  The fixed reflexivity
clause is a theorem.  The soundness of the new Open rule was proved directly
using these readings.

The separation axiom in the ordinary realization must be stated.  Positive
computation equations alone also hold in a singleton structure.  Native
falsity being encoded by $0=S0$ does not justify silently removing the
ordinary falsity bridge.

\subsection{All ordinary logical operations are available}

Put $q_yA=D_{w,y}A$ for an unused $w$.  Define the selector on equations by
\[
 \operatorname{Sel}(v,s=t,r=u)
  := C(v,s,r)=C(v,t,u),
\]
and recursively on equally shaped paired formulas by moving the selector
into their matrices, freshening bound variables.  To combine formulas of
different shapes, pad an equation $E$ by the provably equivalent compound
envelope $q_z\operatorname{Gate}(z,E)$ and then pad the shorter binder
depth.  Here
\[
 \operatorname{Gate}(v,s=t):=C(v,0,s)=C(v,0,t),
\]
and gating also moves recursively through paired binders.  Thus the zero
gate is provable and its successor instances are equivalent to the input.
The inherited native logical interface gives the following macros:
\[
\begin{aligned}
 \bot^* &:=0=S0,\\
 A\to^*B &:=q_y\operatorname{Sel}(y,A,B),\\
 A\wedge^*B &:=q_y\operatorname{Gate}
           (y,\operatorname{Sel}(P(y),A,B)),\\
 \exists^*x A(x)&:=D_{x,y}\operatorname{Gate}(y,A(x)),\\
 A\vee^*B&:=\exists^*v\operatorname{Sel}(v,A,B),\\
 \forall^*x A(x)&:=A(0)\wedge^*q_xA(x).
\end{aligned}
\]
All auxiliary variables are fresh.  The padding is a syntactic
construction, not a new axiom.  Clause use and the recovered ordinary
paired elimination give closure specialization.  Numerical cases give
the selector endpoints; the compound envelope gives unrestricted case
introduction.  These recover every intuitionistic logical rule for the
displayed macros.  No induction on an arbitrary quantified formula is
introduced by this logical recovery.

For example, under $A$ a proof of $B$ gives a proof of the successor
instance of $\operatorname{Sel}(y,A,B)$ from its zero instance.
Closure introduction gives $A\to^*B$.  Closure specialization at $1$
gives its elimination rule.  For an existential, the gated zero instance
is provable and the successor instance is the witnessed formula.
For a conjunction, specialization at $1,2$ selects its two components
using the two parity endpoints.

\subsection{The numerical interpretations}

An iterator translates to the primitive recursive function defined by
iteration with an update that ignores the recursion counter.  iPRA proves
its equations, equation induction, coverage, and separation.

In the opposite direction the inherited scalar development proves, in
this order, a predecessor, reflection of actual equality by a numerical
test, characteristic equations for all quantifier-free formulas, and
pairing with projections.  It is important to preserve that order:
neither pairing nor a faithful equality detector is assumed in obtaining
the predecessor.

If $\pi,L,R$ are the resulting pairing and projection terms, then a
primitive recursion
\[
 f(0,\vec p)=g(\vec p),\qquad
 f(Sn,\vec p)=h(n,f(n,\vec p),\vec p)
\]
is translated by
\[
\begin{aligned}
 V(n,\vec p)&=
 I\bigl(n,\pi(0,g^*(\vec p));
      u.\pi(SL(u),h^*(L(u),R(u),\vec p))\bigr),\\
 f^*(n,\vec p)&=R(V(n,\vec p)).
\end{aligned}
\]
Equation induction proves $L(V(n,\vec p))=n$, and the two projection
identities prove the defining equations for $f^*$.  The construction
recurses externally over each fixed finite primitive recursive
description.  It does not quantify over arbitrary functions internally.

\subsection{Whole domains and return translations}

The interpretation domain in both directions is defined by $x=x$.
Interpreted equality is actual equality.  Uniqueness of a fixed iteration
follows by equation induction on the equality of two solutions with the
same seed and update.  The analogous argument applies to primitive
recursion.  External induction on each finite function description then
proves that both numerical term composites equal the original term on
the entire numerical domain.  Both numerical composite isomorphisms have
graph $x=y$, including at nonstandard elements.

The logical macros and the ordinary reading have formula round trips by
structural induction.  They cover arbitrarily nested quantifiers and
arbitrary retained contexts.  The new opening calculus is equivalent to
the canonical reference on the entire clause consequence relation, so
this existing bridge transfers without changing a domain or an equality.
In particular, the unchanged contradiction $0=S0$ gives equiconsistency.

These are the inherited bridge results formalized in the pinned
\texttt{BiInterpretation}, \texttt{NativeClausal}, and
\texttt{NativeBiInterpretation} modules.\cite{native} Relevant endpoints include
\texttt{biInterp}, \texttt{cprf\_iff\_nprf},
\texttt{native\_iff\_pra}, \texttt{pra\_iff\_native},
\texttt{native\_roundTrip}, \texttt{pra\_native\_roundTrip}, and
\texttt{native\_inconsistent\_iff}.  The new deterministic-opening
comparison and the full-normalization presentation remain written proofs;
they are not additional claimed Lean endpoints.

